% Options for packages loaded elsewhere
\PassOptionsToPackage{unicode}{hyperref}
\PassOptionsToPackage{hyphens}{url}
\PassOptionsToPackage{dvipsnames,svgnames,x11names}{xcolor}
%
\documentclass[
  10pt,
  a4paper,
]{article}
\usepackage{amsmath,amssymb}
\usepackage{setspace}
\usepackage{iftex}
\ifPDFTeX
  \usepackage[T1]{fontenc}
  \usepackage[utf8]{inputenc}
  \usepackage{textcomp} % provide euro and other symbols
\else % if luatex or xetex
  \usepackage{unicode-math} % this also loads fontspec
  \defaultfontfeatures{Scale=MatchLowercase}
  \defaultfontfeatures[\rmfamily]{Ligatures=TeX,Scale=1}
\fi
\usepackage{lmodern}
\ifPDFTeX\else
  % xetex/luatex font selection
  \setmainfont[]{Liberation Serif}
  \setsansfont[]{Liberation Sans}
  \setmonofont[]{DejaVu Sans Mono}
\fi
% Use upquote if available, for straight quotes in verbatim environments
\IfFileExists{upquote.sty}{\usepackage{upquote}}{}
\IfFileExists{microtype.sty}{% use microtype if available
  \usepackage[]{microtype}
  \UseMicrotypeSet[protrusion]{basicmath} % disable protrusion for tt fonts
}{}
\makeatletter
\@ifundefined{KOMAClassName}{% if non-KOMA class
  \IfFileExists{parskip.sty}{%
    \usepackage{parskip}
  }{% else
    \setlength{\parindent}{0pt}
    \setlength{\parskip}{6pt plus 2pt minus 1pt}}
}{% if KOMA class
  \KOMAoptions{parskip=half}}
\makeatother
\usepackage{xcolor}
\usepackage[margin=22mm]{geometry}
\setlength{\emergencystretch}{3em} % prevent overfull lines
\providecommand{\tightlist}{%
  \setlength{\itemsep}{0pt}\setlength{\parskip}{0pt}}
\setcounter{secnumdepth}{-\maxdimen} % remove section numbering
\ifLuaTeX
\usepackage[bidi=basic]{babel}
\else
\usepackage[bidi=default]{babel}
\fi
\babelprovide[main,import]{korean}
\ifPDFTeX
\else
\babelfont{rm}[]{Liberation Serif}
\fi
% get rid of language-specific shorthands (see #6817):
\let\LanguageShortHands\languageshorthands
\def\languageshorthands#1{}
\usepackage{xeCJK}
\setCJKmainfont{Noto Serif CJK KR}
\setCJKsansfont{Noto Sans CJK KR}
\setCJKmonofont{Noto Sans Mono CJK KR}
\xeCJKsetup{CJKspace=true}
\usepackage{xurl}
\usepackage{fvextra}
\usepackage{fancyhdr}
\AtBeginDocument{\DefineVerbatimEnvironment{verbatim}{Verbatim}{breaklines=true,breakanywhere=true,fontsize=\small}}
\setlength{\emergencystretch}{4em}
\setlength{\headheight}{15pt}
\pagestyle{fancy}\fancyhf{}
\fancyhead[L]{\small 중성자별 구조와 비방사 진동}
\fancyfoot[C]{\thepage}
\renewcommand{\headrulewidth}{0.3pt}
\setcounter{tocdepth}{2}
\ifLuaTeX
  \usepackage{selnolig}  % disable illegal ligatures
\fi
\usepackage{bookmark}
\IfFileExists{xurl.sty}{\usepackage{xurl}}{} % add URL line breaks if available
\urlstyle{same}
\hypersetup{
  pdftitle={A16 · npe beta equilibrium과 압축응답},
  pdflang={ko-KR},
  colorlinks=true,
  linkcolor={Maroon},
  filecolor={Maroon},
  citecolor={Blue},
  urlcolor={Blue},
  pdfcreator={LaTeX via pandoc}}

\title{A16 · npe beta equilibrium과 압축응답}
\usepackage{etoolbox}
\makeatletter
\providecommand{\subtitle}[1]{% add subtitle to \maketitle
  \apptocmd{\@title}{\par {\large #1 \par}}{}{}
}
\makeatother
\subtitle{과제 큐시트}
\author{}
\date{}

\begin{document}
\maketitle

\setstretch{1.12}
\textbf{선행:} A15, A02　 \textbf{독립 작업:} 4--7시간 (계획치)

\subsection{밀도를 바꿀 때 조성도 다시 풀어야
하는가?}\label{uxbc00uxb3c4uxb97c-uxbc14uxafc0-uxb54c-uxc870uxc131uxb3c4-uxb2e4uxc2dc-uxd480uxc5b4uxc57c-uxd558uxb294uxac00}

한 밀도에서 beta equilibrium을 구했다. 이 조성을 고정한 미분과 새
밀도에서 평형을 다시 구한 미분이 같은 질문인지 먼저 설명하자.

먼저 한두 문장으로 답한 뒤 아래 설명을 읽자. 첫 예측이 틀렸거나 아직
모르겠다고 적는 것은 감점하지 않는다.

\subsection{필요한 관계와 읽을
위치}\label{uxd544uxc694uxd55c-uxad00uxacc4uxc640-uxc77duxc744-uxc704uxce58}

\[
F(x)=4S(n)(1-2x)-\mu_e(nx),\qquad 0\leq x\leq\frac12.
\] 두 에너지 항을 같은 단위로 평가한다. \(n_e=n_p=nx\)이고
\(\alpha=1-2x\)다.

\begin{itemize}
\tightlist
\item
  \path|learner/src/eos/ns_beta_equilibrium.f90|
\item
  \path|notes/worked/nuclear_calibration.md|
\end{itemize}

\subsection{계산으로
옮기기}\label{uxacc4uxc0b0uxc73cuxb85c-uxc62euxae30uxae30}

기본 경로는 제공 함수의 호출과 해석이다. 아래 \texttt{learner} 시험은
같은 물리 함수를 직접 작성한 경우에 실행한다. 제공 예제를 검사할 때에는
명시적으로 \texttt{calculation}을 선택한다.

\begin{enumerate}
\def\labelenumi{\arabic{enumi}.}
\tightlist
\item
  common nucleon rest mass와 electron total energy의 convention을
  적는다.
\item
  {[}0,1/2{]}에서 residual 양끝과 단조성을 확인한다.
\item
  electron helper를 호출하며 x\_p를 찾는다.
\item
  n±hn에서 조성을 다시 풀어 압력 first law와 equilibrium χ를 검사한다.
\end{enumerate}

\newpage

\subsection{실행해서 무엇을 볼
것인가}\label{uxc2e4uxd589uxd574uxc11c-uxbb34uxc5c7uxc744-uxbcfc-uxac83uxc778uxac00}

패키지 루트에서 실행한다. \texttt{-\/-plan}은 경로만 확인하며 계산
성공을 뜻하지 않는다.

\begin{verbatim}
python tools/run_assignment.py A16 --plan
python tools/run_assignment.py A16 --run \
  --workspace learner --output ../A16_result.json
\end{verbatim}

\subsection{결과를
설명하기}\label{uxacb0uxacfcuxb97c-uxc124uxba85uxd558uxae30}

전하중성, 화학평형, 압축응답은 연결되지만 같지 않다. 두 미분을 모두
계산했다면 어느 쪽을 Cowling 입력으로 선택했는지 명시한다.

\textbf{남길 증거.} 한 밀도의 x\_p,α,μ\_e,잔차; 전하중성 검사; 재평형
미분과 고정조성 미분의 차이

\textbf{판단 기준.} 2점 예: chemical residual만이 아니라 energy-pressure
consistency도 별도 차분으로 제시한다. 최초 예측의 정오보다, 결과를 본 뒤
이유를 고치고 다른 입력에도 적용하는지를 본다.

\subsection{막혔을 때}\label{uxb9c9uxd614uxc744-uxb54c}

\textbf{힌트 1.} S\textgreater0에서 residual의 x\_p 미분 부호를 먼저
생각한다.

\textbf{힌트 2.} 압력 미분을 할 때 주변 밀도의 beta equilibrium을 다시
풀었는지 확인한다.

\subsection{다음 만남에 다시
묻기}\label{uxb2e4uxc74c-uxb9ccuxb0a8uxc5d0-uxb2e4uxc2dc-uxbb3buxae30}

다음 회차에는 전자 화학퍼텐셜과 핵물질 항의 단위를 숨긴 입력을 읽는다.
근 찾기 전에 에너지 convention과 조성 제약을 복원한다. 먼저 짧게 답한 뒤
본문과 자신의 기록을 열어 확인한다.

\textbf{범위.} npe-only toy는 muon과 reaction-rate dynamics를 포함하지
않는다. 조성 평형이 진동 시간척도에서 유지된다는 별도 주장을 하지
않는다.

\end{document}
